\chapter{Sodium (Blood)}

\section*{What Is This Test?}
The sodium blood test measures the concentration of sodium ions (Na\textsuperscript{+}) in the blood serum or plasma. Sodium is the primary extracellular cation and is critical for maintaining fluid balance, nerve impulse transmission, and muscle contraction. Sodium concentration reflects the balance between total body sodium and total body water, making it fundamentally a measure of water homeostasis rather than sodium balance alone. Approximately 85\% of body sodium is in the extracellular fluid compartment, with the remaining 15\% intracellular. The body tightly regulates serum sodium levels through antidiuretic hormone (ADH/vasopressin), thirst, and the renin-angiotensin-aldosterone system. Normal sodium is maintained between 135--145 mEq/L. The test is part of the basic metabolic panel (BMP) and comprehensive metabolic panel (CMP), and is one of the most frequently ordered laboratory tests worldwide.

\section*{Alternative Names}
Na; Sodium, Serum; Serum Sodium; Plasma Sodium; Blood Sodium; Na+ Level; Electrolyte Panel Component

\section*{Principle}
Serum sodium is measured using ion-selective electrode (ISE) potentiometry. In direct ISE, the sample is introduced undiluted and the sodium-sensitive glass membrane electrode develops a potential proportional to the logarithm of sodium ion activity in the sample. In indirect ISE, the sample is first diluted with a buffer before measurement. Direct ISE (used in blood gas analyzers and point-of-care devices) measures sodium activity in the plasma water phase and is not affected by pseudo-hyponatremia from hyperlipidemia or hyperproteinemia. Indirect ISE (used in most automated chemistry analyzers like Roche Cobas, Abbott Architect, Siemens Atellica) measures sodium concentration in total plasma volume and therefore can be falsely low when the plasma water fraction is reduced by very high lipid or protein concentrations. The sodium-selective electrode typically uses a glass membrane containing aluminum silicate (NAS 11-18 glass). The reference electrode completes the electrochemical circuit. Modern analyzers incorporate built-in quality control and can results in under 30 seconds.

\section*{History}
The role of sodium in physiology was first explored by Claude Bernard in the 1850s. James Gamble's seminal work in the 1940s characterized electrolyte composition of body fluids. Flame photometry was introduced for sodium measurement in the 1940s, using the characteristic yellow-orange wavelength (589 nm) emitted by sodium atoms excited in a flame. The first ion-selective electrode for sodium was developed in the 1960s, replacing flame photometry in most clinical laboratories by the 1980s. The glass electrode for sodium measurement was pioneered by George Eisenman at the University of Chicago in the 1960s. Automated multichannel analyzers (Technicon SMA 6/60 in the 1960s) made electrolyte panels routine clinical practice.

\section*{Why This Test Is Ordered}
\begin{itemize}
  \item Evaluation of hyponatremia (confusion, seizures, weakness, nausea) or hypernatremia (thirst, altered mental status, muscle twitching)
  \item Monitoring of fluid and electrolyte status in patients receiving IV fluids, diuretics, or with known renal, cardiac, or hepatic impairment
  \item Part of routine electrolyte panel monitoring in patients taking medications that affect sodium balance: thiazide diuretics, loop diuretics, SSRIs, carbamazepine, cyclophosphamide, NSAIDs, desmopressin, amiodarone
  \item Evaluation of acid-base disorders; serum sodium contributes to calculation of serum osmolality: 2(Na) + Glucose/18 + BUN/2.8
  \item Workup of adrenal insufficiency, SIADH, diabetes insipidus, and cerebral salt wasting
  \item Preoperative assessment and monitoring during major surgery
  \item Monitoring refeeding syndrome in malnourished patients
\end{itemize}

\section*{Is This a Primary or Secondary Test}
Sodium measurement is a primary test for evaluation of fluid and electrolyte status. It is a component of the basic metabolic panel (BMP) and comprehensive metabolic panel (CMP), which are first-line screening tests in many clinical settings. Sodium results are always interpreted together with clinical volume status assessment and plasma osmolality for accurate diagnosis.

\section*{How This Test Is Performed}
A healthcare professional will draw blood from a vein in your arm using a small needle. A tourniquet is applied briefly to locate the vein. After the needle is inserted, blood is collected into a serum separator tube (gold-top or red-top) or a lithium heparin tube (green-top). The tourniquet should be released as soon as blood begins flowing to avoid hemoconcentration, which can falsely elevate sodium. The sample is centrifuged and sodium is measured on the chemistry analyzer. Results are typically available within 30 minutes to 1 hour. For patients with difficult venous access, a heel stick (infants) or finger stick capillary sample can be used with point-of-care devices, though accuracy may be lower.

\section*{What Preparation Is Needed}
No specific patient preparation is required for a serum sodium test. Fasting is not necessary. However, avoid excessive water intake immediately before the test as it can cause dilutional changes. The healthcare provider should be informed of all medications, particularly diuretics, lithium, desmopressin, oxcarbazepine, carbamazepine, SSRIs (especially in elderly patients), and IV fluids currently being administered. Recent IV saline infusion can transiently elevate sodium. Prolonged tourniquet application during phlebotomy can cause hemoconcentration and falsely elevate sodium by 1--2 mEq/L. The sample should be processed within 4 hours if stored at room temperature or within 24 hours if refrigerated at 2--8\textdegree C.

\section*{Contraindications and Precautions}
No absolute contraindications exist for venipuncture. Pre-analytical factors that can affect sodium results: hemolysis (potassium is released from cells but sodium effect is minimal), prolonged tourniquet application (false increase), drawing blood proximal to an IV infusion site (massive false change), and using the wrong tube additive. Pseudohyponatremia occurs with indirect ISE methods when plasma is significantly lipemic (triglycerides >1500 mg/dL) or when total protein exceeds 10 g/dL, as the lipid/protein displaces plasma water, reducing apparent sodium per total volume. This is not seen with direct ISE. Hyperglycemia causes true dilutional hyponatremia via osmotic water shift from cells: for every 100 mg/dL increase in glucose above 100 mg/dL, serum sodium decreases by approximately 1.6--2.4 mEq/L. Mannitol and sorbitol infusions can cause similar shifts. Diuretic use, recent surgery (increased ADH), and conditions affecting ADH secretion (pain, nausea, stress) must be considered in interpretation.

\section*{Risks}
Minimal risk associated with venipuncture. Mild pain or bruising at the puncture site is common and resolves within days. Rare complications include hematoma formation, infection at the site (extremely rare), nerve injury if the needle contacts a superficial nerve, and vasovagal syncope (fainting), particularly in anxious patients. No risk of significant blood loss; typical blood volume drawn is 3--5 mL.

\section*{What Happens After the Test}
A bandage is applied to the puncture site and should remain for 15--30 minutes. Results are typically available within 30 minutes to 1 hour. Critical values: sodium <120 mEq/L or >160 mEq/L require immediate clinical intervention due to risk of seizures, cerebral edema, or permanent neurological injury. The rate of sodium correction is critical: chronic hyponatremia (>48 hours duration) must be corrected slowly at no more than 8--10 mEq/L per 24 hours to avoid osmotic demyelination syndrome (central pontine myelinolysis). Acute hyponatremia (<48 hours) can be corrected more rapidly. Hypernatremia correction should not exceed 10--12 mEq/L per 24 hours to avoid cerebral edema. Following abnormal results, your healthcare provider will assess volume status (hypovolemic, euvolemic, or hypervolemic) to determine the underlying etiology.

\section*{Understanding Results}

\subsection*{Below Normal}
\begin{itemize}
  \item \textbf{Hyponatremia (Na <135 mEq/L):} The most common electrolyte abnormality in hospitalized patients, occurring in 15--30\% of inpatients.
  \item \textbf{Hypovolemic hyponatremia:} Sodium and water are both depleted, but sodium loss exceeds water loss. Caused by GI losses (vomiting, diarrhea, nasogastric suction), renal losses (diuretic excess, salt-wasting nephropathy, mineralocorticoid deficiency in Addison's disease, cerebral salt-wasting syndrome), third-space losses (burns, pancreatitis), and excessive sweating. Serum osmolality is low (<275 mOsm/kg) and urine sodium is low (<20 mEq/L in extrarenal losses, >20 mEq/L in renal losses).
  \item \textbf{Euvolemic hyponatremia:} Most commonly caused by SIADH (syndrome of inappropriate antidiuretic hormone secretion), which can result from CNS disorders (meningitis, encephalitis, stroke, tumor, trauma), pulmonary disorders (pneumonia, TB, small cell lung cancer), medications (SSRIs, carbamazepine, cyclophosphamide, NSAIDs), pain, nausea, or prolonged surgery. Also caused by psychogenic polydipsia (water intake >18--20 L/day), hypothyroidism, adrenal insufficiency, and reset osmostat. Urine osmolality is inappropriately >100 mOsm/kg with urine sodium >20--40 mEq/L.
  \item \textbf{Hypervolemic hyponatremia:} Total body sodium and water are both increased, but water increase exceeds sodium. Caused by congestive heart failure, cirrhosis with ascites, nephrotic syndrome, and advanced chronic kidney disease. Urine sodium may be low (<20 mEq/L) in heart failure and cirrhosis due to effective arterial underfilling.
  \item \textbf{Sodium <125 mEq/L:} Risk of cerebral edema, headache, nausea, confusion. <120 mEq/L: seizures, coma, respiratory arrest risk.
\end{itemize}

\subsection*{Above Normal}
\begin{itemize}
  \item \textbf{Hypernatremia (Na >145 mEq/L):} Always reflects hyperosmolality (>295 mOsm/kg) and indicates a relative deficit of water in relation to total body sodium. Seen in 1--3\% of hospitalized patients; mortality approaches 40--60\% when Na >160 mEq/L in hospitalized adults.
  \item \textbf{Hypovolemic hypernatremia (low total body sodium):} Extrarenal causes: diarrhea (especially osmotic diarrhea from lactulose), vomiting, excessive sweating, burns. Renal causes: loop diuretics, osmotic diuresis (hyperglycemia, mannitol), post-obstructive diuresis, recovery from ATN. Urine sodium >20 mEq/L in renal losses; <10 mEq/L in extrarenal losses.
  \item \textbf{Euvolemic hypernatremia (normal total body sodium):} Diabetes insipidus (central or nephrogenic), insensible water losses (fever, hyperventilation, exposure to high environmental temperatures). Central DI: lack of ADH production (pituitary surgery, head trauma, tumors, sarcoidosis, histiocytosis). Nephrogenic DI: renal resistance to ADH (lithium therapy, hypercalcemia, hypokalemia, chronic kidney disease, genetic mutations in aquaporin-2 or V2 receptor). Urine osmolality low (<300 mOsm/kg) with polyuria.
  \item \textbf{Hypervolemic hypernatremia (high total body sodium):} Iatrogenic administration of hypertonic saline, sodium bicarbonate infusion, mineralocorticoid excess (primary hyperaldosteronism, Cushing's syndrome), seawater ingestion.
  \item \textbf{Sodium >155 mEq/L:} Neurological symptoms prominent---irritability, hyperreflexia, seizures. >160 mEq/L: altered mental status, coma, irreversible neurological damage.
\end{itemize}

\section*{Normal Results}
\begin{itemize}
  \item \textbf{Adults:} 135--145 mEq/L (mmol/L)
  \item \textbf{Children (1--16 years):} 135--145 mEq/L
  \item \textbf{Infants (0--1 year):} 134--146 mEq/L
  \item \textbf{Newborns (0--7 days):} 133--146 mEq/L
  \item \textbf{Plasma osmolality (calculated):} 2(Na) + Glucose(mg/dL)/18 + BUN(mg/dL)/2.8 = 275--295 mOsm/kg
  \item \textbf{Critical low value:} <120 mEq/L
  \item \textbf{Critical high value:} >160 mEq/L
\end{itemize}

\section*{Follow-up Tests}
\begin{itemize}
  \item \textbf{Plasma osmolality (measured):} To confirm true hyponatremia vs. pseudohyponatremia or hypertonic hyponatremia (from hyperglycemia or mannitol). Measured by freezing point depression osmometry.
  \item \textbf{Urine osmolality:} To differentiate SIADH (urine osmolality >100 mOsm/kg, often > plasma) from primary polydipsia (urine maximally dilute, <100 mOsm/kg).
  \item \textbf{Urine sodium concentration:} <20 mEq/L suggests effective circulating volume depletion (heart failure, cirrhosis, dehydration); >20--40 mEq/L suggests SIADH, renal salt wasting, diuretic use, or adrenal insufficiency.
  \item \textbf{Comprehensive metabolic panel (CMP):} To assess potassium (often low in diuretic-induced hyponatremia), BUN/creatinine (renal function), glucose (hyperglycemic hyponatremia), and albumin (liver disease).
  \item \textbf{Serum uric acid:} Low in SIADH (<4 mg/dL) due to dilution and increased clearance; normal or elevated in volume depletion.
  \item \textbf{Thyroid function tests (TSH, free T4):} Hypothyroidism causes euvolemic hyponatremia.
  \item \textbf{Morning cortisol / ACTH stimulation test:} To rule out adrenal insufficiency as cause of hyponatremia.
  \item \textbf{Water deprivation test with desmopressin challenge:} To differentiate central vs. nephrogenic diabetes insipidus in hypernatremia workup.
  \item \textbf{Brain imaging (CT or MRI):} If SIADH is suspected from CNS pathology or if osmotic demyelination syndrome is a concern.
\end{itemize}

\section*{References}
\begin{enumerate}
  \item Burtis CA, Ashwood ER, Bruns DE. Tietz Textbook of Clinical Chemistry and Molecular Diagnostics. 7th ed. Elsevier; 2023.
  \item Wallach J. Interpretation of Diagnostic Tests. 10th ed. Wolters Kluwer; 2022.
  \item Tierney LM, Saint S, Drazen JM. CURRENT Medical Diagnosis and Treatment. McGraw-Hill; 2023.
  \item Fischbach FT, Dunning MB. A Manual of Laboratory and Diagnostic Tests. 10th ed. Wolters Kluwer; 2023.
  \item Spasovski G, Vanholder R, Allolio B, et al. Clinical practice guideline on diagnosis and treatment of hyponatraemia. Intensive Care Med. 2014;40(3):320--331.
  \item Verbalis JG, Goldsmith SR, Greenberg A, et al. Diagnosis, evaluation, and treatment of hyponatremia: expert panel recommendations. Am J Med. 2013;126(10 Suppl 1):S1--S42.
  \item Adrogu\'{e} HJ, Madias NE. Hypernatremia. N Engl J Med. 2000;342(20):1493--1499.
  \item Sterns RH. Disorders of plasma sodium---causes, consequences, and correction. N Engl J Med. 2015;372(1):55--65.
\end{enumerate}

\clearpage
